% Transcription par lecture visuelle directe, sans OCR.
% Statut : transcrit ; relecture indépendante à effectuer.
% Source : PDF original, vue 148, côté droit, page manuscrite 297, après le chapitre treizième.
\chapter{Préfecture de M. Matthieu Garnier. Établissement des Jésuites et de la Congrégation de Paris. M. Recorbet nommé père de la Congrégation.}

\section*{M. Garnier préfet}
% Mention marginale de la page manuscrite 297.
La Congrégation après avoir été agitée en diverses manières par les vicissitudes, à travers lesquelles il avait plu à la Providence de la conduire, était enfin rentrée dans le calme le plus parfait ; mais les devoirs que cette position nouvelle imposait au Préfet, n'en demandaient pas moins toute son attention. S'il n'était plus question maintenant de lutter corps à corps contre des difficultés sans cesse renaissantes qui venaient menacer la société jusque dans son existence, il fallait travailler à la consolider, saisir avec habileté toutes les occasions d'exciter, de ranimer et de réchauffer le zèle des Congréganistes. Il fallait surtout faciliter le développement du bien qu'on pouvait attendre de leurs efforts réunis. Notre institution mûrie par les tribulations et par l'expérience de plusieurs années, ne pouvait plus être considérée comme une institution nouvelle, il était temps qu'on en vît enfin porter des fruits abondants.

Lorsque la volonté du Souverain maître n'appelle pas la Congrégation d'une manière spéciale à faire arriver son règne, ce qu'elle prend alors soin de lui faire connaître
% Source : PDF original, vue 149, côté gauche, page manuscrite 298.
par les circonstances où elle la place, il lui reste toujours le moyen le plus ordinaire, le moyen de tous les moments de remplir son but principal qui est de procurer la gloire de Dieu ; il lui reste, disons nous, le concours de ses sections. C'est donc là ce qui réclame par dessus tout tous les soins du préfet. Il doit faire comprendre de plus en plus aux membres de ces sections que le but de leurs œuvres est beaucoup moins de répandre des bienfaits que de sauver des âmes ; il doit engager par tous les moyens qui sont en son pouvoir, les Congréganistes, à apporter dans ces œuvres cet esprit de zèle, d'humilité et d'abnégation de sa propre volonté qui peut faire espérer de se sanctifier soi-même en sanctifiant les autres.

\section*{Travaux des sections}
% Mention marginale de la page manuscrite 298.
M. Garnier comprit toute la nécessité qu'il y avait de ne pas négliger cette partie de son administration. Sous son gouvernement, non seulement on vit les sections déployer de nouveau toute l'activité dont elles avaient fait preuve dans les meilleurs temps de la Congrégation mais encore se faire remarquer par de nouveaux progrès.

La section des aumônes avait repris dès que cela avait été possible le cours complet de ses visites de charité. Elle augmenta le nombre de ses pauvres, et plus que jamais elle s'efforça de joindre l'aumône spirituelle à l'aumône temporelle.

La section de l'instruction eut promptement rétabli les catéchismes qu'elle faisait dans les paroisses ; elle en forma de nouveaux. Les uns et les autres devinrent de plus en plus nombreux. Nous aurons occasion de reparler plus tard de cette œuvre ; nous la verrons alors prendre plus d'extension
% Source : PDF original, vue 149, côté droit, page manuscrite 299.
et surtout plus de fixité.

Dans le cours de 1816, la section de l'instruction établit à St François un catéchisme pour les militaires. Il est à regretter que cette institution n'ait pas pu se soutenir bien longtemps. Ce fut à l'intervention de quelques officiers de la garnison, fervents catholiques, que nous dûmes la formation de ce catéchisme. Les Congréganistes de l'époque n'ont pas oublié la piété de M. de Chateaubriant, neveu du célèbre écrivain, celle du capitaine Martin, du brave lieutenant Bouquet, et de plusieurs autres officiers, comme eux chevaliers sans peur et sans reproche. Souvent nous nous rencontrions à Fourvières ; nous entourions ensemble l'autel de Marie ; nous participions ensemble au divin banquet. En un mot ces défenseurs de l'autel et du trône avaient fini par contracter avec plusieurs d'entre nous auprès de notre bonne mère une union qui les rendait presque Congréganistes. Ils cherchèrent les moyens de faire instruire les militaires qui leur étaient subordonnés, et grâce à notre section de l'instruction, nous pûmes les aider dans cette bonne œuvre. Il était touchant de voir 50 ou 60 soldats, quelquefois davantage, prêter une oreille attentive aux explications des vérités de la religion, explications qui sortaient bien souvent de la bouche d'un tout jeune homme. On les entendait ensuite chanter à gorge déployée des cantiques composés exprès pour eux et imprimés aux frais de la section. C'était un moyen de placer sur leurs lèvres et de faire parvenir jusqu'à leurs cœurs, des sentiments de foi et de dévouement à la cause de la religion, exprimés
% Source : PDF original, vue 150, côté gauche, page manuscrite 300.
avec ce ton de franchise et de simplicité qui convient à des militaires.

Ce fut vers le même temps que prit naissance une des œuvres les plus intéressantes de la section de l'instruction, celle des Savoyards. Le vénérable M. Cattet aujourd'hui curé de St Paul était alors simplement attaché à cette paroisse. Il consacrait la plus grande partie de son temps à l'exercice des bonnes œuvres ; à l'exemple de plusieurs Saints prêtres et notamment des jésuites, il s'était dévoué surtout à la sanctification des âmes dans les rangs les plus obscurs et les plus humbles de la société. La classe qui par dessus tout avait excité son zèle, était celle de ces enfants si nombreux qui abondent dans notre ville pour y exercer la profession de décrotteurs ou autres professions analogues. La police municipale était alors animée des meilleures intentions ; elle ne pouvait qu'applaudir au zèle de ce vertueux ecclésiastique, et reconnaître l'immense avantage de son intervention en faveur d'une jeunesse que l'absence de principes et d'éducation pouvait conduire à tous les désordres. Elle avait donc accepté pour elle le patronage que M. Cattet avait bien voulu lui offrir, et en les plaçant sous son autorité, elle avait soumis à certaines règles tous ceux qui à l'avenir voudraient exercer à Lyon l’état de décrotteur. Nous avons à déplorer que ces règlements de police n'aient pas été maintenus pendant bien longtemps.

M. Cattet eut recours à la Congrégation pour le seconder. La section de l'instruction fut donc mise en possession de l'œuvre. Bientôt elle y joignit le soin des ramoneurs et de
% Source : PDF original, vue 150, côté droit, page manuscrite 301.
tous les autres industriels de bas étage qui viennent à Lyon chaque année chercher des moyens d'existence. Elle trouva là un ample aliment à son zèle. Grâce à cette bonne œuvre un grand nombre de ces pauvres jeunes gens trouve les moyens de persévérer dans les principes de religion qu'ils ont reçus dans leurs villages. Leurs curés respectifs ne manquent pas, lorsqu'ils se rendent à Lyon, de les recommander à M. Cattet, et celui-ci veille sur eux avec toute la sollicitude d'un père.

La section des prisons et hôpitaux commença à la même époque à donner aux visites qu'elle fait aux malades de l'hôtel-dieu, ce grand développement que cette œuvre a conservé depuis. Les obstacles qui l'avaient écartée des prisons civiles, même avant le commencement de la persécution subsistaient encore ; ce ne fut qu'un peu plus tard qu'on parvint à les aplanir ; mais les membres de la section avaient trouvé un dédommagement dans l'œuvre de la prison militaire des recluses ; cette œuvre était pour eux une source de consolation. Les pauvres soldats détenus recevaient avec attention et plaisir les instructions qui leur étaient adressées. Souvent le retour de plusieurs d'entre eux aux maximes et aux pratiques de la religion devenait un grand sujet de joie qui récompensait amplement de tous leurs soins ceux qui étaient chargés de cette bonne œuvre.

La bonté infinie de notre Dieu tient toujours en réserve pour les faire goûter à la Congrégation quelques consolations nouvelles. Celle qu'il lui ménagea dans le cours de cette année 1816 y produisit une impression bien vive et bien durable. Vous vous souvenez sans doute, que dans l'année 1804, notre société
% Source : PDF original, vue 151, côté gauche, page manuscrite 302.
avait contracté une alliance avec une autre société pieuse établie à Chambéry, sous le nom des amis de Chambéry, et dont M. Rey, aujourd'hui Mgr l'évêque d'Annecy, était le directeur.

\section*{Visite de M. Rey}
% Mention marginale de la page manuscrite 302.
Quoique cette société ne subsistât plus les relations établies entre M. l'abbé Rey et la Congrégation de Lyon, n'avaient pas à cause de cela cessé d'exister. Bien au contraire elles étaient peut-être devenues plus intimes, car ainsi que nous avons eu occasion de vous le dire, M. Rey avait souffert en même temps que nos amis et pour la même cause ; il avait été prisonnier par suite de sa coopération à la transmission de la correspondance du Pape. Or les persécutions supportées ainsi en commun pour la justice établissent entre les chrétiens une véritable confraternité ; toutefois l'occasion ne s'était pas encore présentée pour nous de connaître, autrement que par sa correspondance ce saint et pieux ecclésiastique. Il était alors devenu vicaire général de Chambéry. Il vint à Lyon en 1816 pour y prêcher le carême dans l'église de St François de Sales.

Dès qu'il fut arrivé plusieurs Congréganistes s'empressèrent de lui rendre visite ; et purent enfin juger de près de toute la chaleur de ce cœur rempli de l'amour de Dieu et du prochain. Presque tous nos frères accoururent aussi souvent que cela leur fut possible à St François de Sales, pour recueillir avec délices ces expressions si expansives qui coulaient comme de leur source, d'une âme où la charité de J. C. surabondait.

La fête de l'Annonciation approchait ; il parut à nos chefs qu'elle arrivait à propos pour nous donner un moyen
% Source : PDF original, vue 151, côté droit, page manuscrite 303.
facile de mettre ce saint prêtre en contact direct avec tous les membres d'une Congrégation qui s'estimait si heureuse de le compter au nombre de ses associés. Il accueillit avec bienveillance la prière qui lui fut adressée de venir en cette solennité célébrer la messe au milieu de nous et nous distribuer en même temps le pain de la parole divine.

Ce fut comme vous pouvez vous le figurer, un véritable jour de bonheur pour la Congrégation que le jour où il lui fut donné de contempler ce digne ministre de J. C. au milieu d'elle. Ce n'était pas seulement par les discours qu'il excitait à la vertu tous ceux qui se trouvaient autour de lui ; tout prêchait en lui pendant la célébration des Saints mystères ; ce recueillement si éminemment religieux, ses gestes et jusqu'au son de sa voix venaient révéler combien il était convaincu de la dignité du sacerdoce et nous pénétrer du respect le plus profond. Ajoutez à cela l'impression que produisaient les paroles brûlantes qu'il jeta parmi nous, comme un torrent de feu, pour embraser nos cœurs de l'ardeur de cette charité qui régnait en souveraine dans le sien, et vous pourrez vous faire peut-être quelque idée des sentiments dont tous les Congréganistes étaient animés en la présence de l'homme de Dieu.

Ce fut alors que la Congrégation contracta avec lui ce pacte d'alliance qui s'exécute si fidèlement de part et d'autre. Il prit envers nous l'engagement de nous mettre tous les jours aux pieds de J. C. notre commun maître en offrant l'adorable sacrifice ; de notre côté nous devions réciter pour lui un \textit{Ave Maria}\footnote{(Note du transcripteur) Littéralement : « Je vous salue, Marie ».} à chacune de nos assemblées.

Vous ne serez pas étonnés que M. Rey nous eût tous
% Source : PDF original, vue 152, côté gauche, page manuscrite 304.
électrisés ; pour juger par vous-même ce que nous devions ressentir alors, vous n'avez qu'à vous rappeler ce que nous éprouvons à chaque troisième dimanche de juillet. Lorsque cette lettre de feu et si bien nommée ce mandement de l'amitié, vient périodiquement chaque année renouveler les mêmes sentiments dans tous les Congréganistes.

M. Rey aimait à dire que le Seigneur a répandu sur certaines âmes une espèce de glu, dont l'effet est de les attacher facilement et si fortement aux âmes de leurs frères, que rien ne peut ensuite rompre l'union qu'elle a cimentée. Il ajoutait, ce dont au reste il n'est guère possible de douter lorsqu'on l'a connu, que le bon Dieu l'avait abondamment pourvu de cette glu, et il s'estimait heureux d'en avoir rencontré beaucoup dans les Congréganistes de Lyon.

\section*{Rétablissement des Jésuites}
% Mention marginale de la page manuscrite 304 ; placée avant le récit de cette restauration.
Pendant que la Congrégation de Lyon, tout en se livrant de nouveau avec zèle à l'exercice des bonnes œuvres, savourait avec délices les petites douceurs dont il plaisait à la Providence de la gratifier, la Congrégation de Paris renaissait de ses cendres plus brillante que jamais. Mais comme cette résurrection se lie à une autre résurrection d'un intérêt bien plus général et dont elle fut la suite et la conséquence, nous nous occuperons d'abord de celle-là, sauf à revenir à l'autre plus tard.

Aussitôt que le Pape Pie VII, de retour à Rome, put enfin donner ses soins aux affaires de l'Église, une de ses premières pensées fut une pensée de douloureux regrets sur la destruction de l'ordre des jésuites. Dieu lui fit comprendre toute l'importance, la presque nécessité du rétablissement de cette troupe auxiliaire si éminemment utile à la cause de la religion et dont la chute,
% Source : PDF original, vue 152, côté droit, page manuscrite 305.
dans le dernier siècle, avait été le prélude de tous les maux qui étaient venus fondre sur l'Église de J. C. Le cœur du père commun des fidèles gémissait sur les plaies profondes que l'irréligion avait faites à la société. Il crut pouvoir espérer que le concours de ces Saints religieux seconderait les efforts qu'il se proposait de tenter pour y apporter un remède efficace. Par une bulle solennelle, il rétablit donc l'ordre des Jésuites ; il appelle au secours de l'Église ces enfants de St Ignace ; il sollicite le concours de ces défenseurs de la foi ; il veut qu'ils viennent se ranger sous ses ordres, combattre avec lui, par la patience, les ennemis de la religion, les vaincre par la douceur, triompher de la dureté de leurs cœurs et rapporter au bercail ces brebis égarées pour lesquelles à la moindre marque de repentir sont prêtes à s'abaisser toutes les barrières qui les séparent encore de ces gras pâturages que le Seigneur lui-même leur a préparés.

Voici encore une occasion, Messieurs, d'admirer comment la Providence sait ménager d'avance à la miséricorde les moyens de réparer les maux occasionnés par les passions humaines qui ont servi d'instrument à la justice. Un peu d'attention s'il vous plaît, nous allons ici la prendre sur le fait.

Pendant que toutes les puissances catholiques, saisies de je ne sais quel esprit de vertige qu'une philosophie mensongère avait fait élever comme un brouillard pour obscurcir tout entendement, sollicitaient avec instance la destruction d'un ordre le plus ferme appui de leurs trônes ; pendant que le Pontife suprême par une foiblesse inexprimable ou un aveuglement incompréhensible, se résignait lui-même à céder à ces demandes insensées et apposait sa signature à un acte
% Source : PDF original, vue 153, côté gauche, page manuscrite 306.
qui renversait d'un trait de plume ce solide rempart qui n'avait pu défendre l'Église et son chef des fureurs de l'impiété, nous voyons apparaître un prince schismatique, l'autocrate de toutes les Russies, qui vient solliciter, exiger presque, la conservation de l'ordre proscrit en faveur de ses états. C'est que Dieu avant la tempête qui allait ébranler le monde, voulait mettre cette précieuse semence à l'abri de l'orage, afin qu'au retour du beau temps elle vînt de nouveau repeupler et embellir le jardin de l'Église dont elle faisait jadis un des plus beaux ornements.

Mais ce n'est pas encore tout ; suivons toujours la Providence. Pie VII élève la voix et rappelle les Jésuites, et au moment précis, l'Empereur de Russie les chasse de ses états. Ce prince schismatique, instrument docile de la volonté du souverain maître, sans toutefois qu'il puisse comprendre que c'est sa main qui le conduit, a conservé et mis à l'abri de tout danger ces saints religieux aussi longtemps que le monde catholique les a repoussés ; aujourd'hui que le monde catholique leur tend les bras et que le chef de l'Église les réclame, la mission que la Providence lui avait donnée est terminée, il se hâte de les renvoyer car le Saint Père les attend ; il faut qu'ils aillent se ranger autour de lui ; ainsi l’a résolu l’arbitre souverain des destinées des peuples. Qui pourrait donc avec un esprit tant soit peu observateur, méconnaître cette action si douce et si forte en même temps, qui sait toujours tout disposer pour qu’au temps marqué les desseins du très haut s'accomplissent.

Pour corroborer ce que nous avançons, et suivre jusque dans ses plus petits détails, la marche que le Seigneur
% Source : PDF original, vue 153, côté droit, page manuscrite 307.
imprime aux événements, ajoutons encore une remarque qui mérite aussi quelque attention. Les Jésuites avaient bien été presque miraculeusement conservés en Russie, mais ils y étaient en petit nombre. Obligés de se disséminer sur tous les points de la chrétienté, s'ils eussent été réduits à leurs propres forces, il se serait écoulé bien des années avant qu'ils eussent pu s'être assez augmentés pour pouvoir rendre à l'Église de nombreux services. Dieu y avait pourvu ; cet ordre des Pères de la Foi, dont nous avons déjà parlé que nous avons aperçu se former en Allemagne et de là s'étendre dans toute l'Europe et notamment en France où malgré leur dispersion, ils sont toujours restés unis de cœur et prêts à se rejoindre au premier signal, cet ordre est là. Il suivait autant que possible la règle des jésuites, il ne s'était établi que dans l'intention bien avouée de faire revivre sous un autre nom les enfants de St Ignace. Cet ordre nouveau, dès le moment que le Pape avait prononcé le rétablissement des Jésuites, devait s’absorber, mais il offrait au petit noyau des Jésuites qui revenait de Russie, comme un immense noviciat dont les nombreux religieux étaient pour la plupart bien disposés et tout préparés à venir grossir leurs rangs. Aussi s'empressa-t-on de les admettre. Les Jésuites en furent notablement accrus et l'ordre des Pères de la Foi se confondit si rapidement avec l'ordre nouvellement rétabli que beaucoup de personnes crurent et qu'il n'est pas rare de rencontrer encore des personnes qui croient que ces deux ordres n'étaient qu'une seule et même chose.

Ceux des jésuites qui vinrent en France s'établirent d'abord à Paris. De là ils s'étendirent sur tous les points de la France et
% Source : PDF original, vue 154, côté gauche, page manuscrite 308.
formèrent en peu d'années de nombreuses maisons d'éducation qui faisaient la joie et la consolation de tous les pères de famille qui désiraient pouvoir donner à leurs enfants une éducation chrétienne. St Acheul, Dole, Forcalquier, Billom et plusieurs autres maisons de ce genre ont laissé des souvenirs qui dureront longtemps. Plaignons le sort de notre malheureux pays dans lequel l'esprit philosophique est encore une fois parvenu à renverser ces pieux établissements.

\section*{Rétablissement de la Congrégation de Paris}
% Mention marginale de la page manuscrite 308.
Ce fut sous les ailes des Jésuites que prit naissance à Paris la nouvelle Congrégation. Ils firent un appel aux principaux d'entre les anciens membres et aussitôt les Montmorency, les Noailles, les Rohan, les Berthier, les Rivière, etc. se hâtèrent d'accourir. Le Père Ronsin, autrefois père de la foi, maintenant devenu Jésuite, fut chargé de la direction de la nouvelle société. Cette nouvelle Congrégation fut formée à peu près sur le plan de l'ancienne, mais loin de prendre quelques mesures pour prévenir les dangers auxquels une trop grande publicité l'avait exposée, on aggrava cette position périlleuse ; car le P. Ronsin qui comme autrefois le P. Delpuy présidait seul aux réceptions, désirant dans l'ardeur de sa charité étendre autant que possible le bien que pouvait procurer aux fidèles ce pieux établissement, ouvrait la porte à deux battants pour laisser entrer tout le monde ; il suffisait d'une simple demande pour obtenir d'en faire partie. Il résulta de là un grand inconvénient qu'il semble qu'on aurait dû prévoir, des personnes se faisaient admettre dans la Congrégation dans l'espérance d'obtenir la protection de ceux des Congréganistes
% Source : PDF original, vue 154, côté droit, page manuscrite 309.
 auxquels le rang qu'ils occupaient dans le monde, ou les places qu'ils remplissaient pouvaient faire supposer quelque influence. On conçoit que des hommes qui ne venaient s'enrôler ainsi par ambition ne pouvaient guère être bons Congréganistes. D'ailleurs cette manière d'agir semblait justifier les attaques dirigées par les ennemis de la religion contre la Congrégation qu'ils accusaient d'avoir des vues politiques, ce qui n'a jamais eu l'ombre de vérité, pas plus pour la Congrégation de Paris que pour celle de Lyon.

La Congrégation de Lyon s'était trop bien trouvée de sa fidélité à suivre ses règles habituelles de prudence pour jamais s'en écarter, aussi pour ne pas sortir de son incognito, se tint-elle sagement à l'écart. Elle laissa bien, à la vérité, ses membres profiter de tous les avantages que pouvait leur offrir la Congrégation lorsqu'ils allaient à Paris, et cela était très facile, car il suffisait de se présenter pour en être de suite mis en possession. Le P. Ronsin, dans sa grande charité, avait pour maxime que tous les Congréganistes de France appartenaient de droit à la Congrégation de Paris, et devaient y avoir leurs entrées ; en conséquence il se hâtait de leur ouvrir toutes les portes.

Mais si la Congrégation de Lyon se montra disposée à accepter de bon cœur tout ce qu'on voulait bien lui offrir, elle se garda bien d'user de réciprocité, c'eût été par trop dangereux. Elle se contenta seulement de recevoir, comme par le passé, ceux des Congréganistes qui avaient appartenu à l'ancienne Congrégation de Paris et qui en cette qualité avaient déjà obtenu leurs entrées dans la nôtre. L'expérience a prouvé combien elle a eu raison de ne pas se laisser entraîner dans un système qui aurait
% Source : PDF original, vue 155, côté gauche, page manuscrite 310.
attiré sur elle tous les dangers qui ont fini par renverser la Congrégation de Paris, et auxquels si elle fût entrée dans cette voie périlleuse, elle n'eût pas mieux échappé.

\section*{Le P. Roger à Paris}
% Mention marginale de la page manuscrite 310.
Parmi les pères de la foi qui vinrent successivement à Paris faire leur noviciat chez les jésuites, pour devenir de véritables enfants de St Ignace, le P. Roger ne fut pas des derniers. Cette circonstance le fixa tout à fait dans cette grande capitale. Il commença dès lors à rétablir au moins par correspondance quelques rapports avec sa chère Congrégation de Lyon. Il fit plus, il devint à Paris le père de tous les Lyonnais, auxquels il était toujours disposé à rendre tous les services spirituels ou même temporels qui dépendaient de lui. Ses entrailles paternelles se dilataient lorsqu'il avait le bonheur de voir un des nôtres. C'était pour lui une jouissance indicible car son cœur était demeuré tout Lyonnais. Si un jésuite qui fait de sa propre volonté une abnégation si entière avait pu conserver un seul désir, c'aurait été certainement pour le P. Roger celui de se retrouver encore avec cette Congrégation de Lyon pour laquelle il avait conservé une si vive tendresse. De son côté la Congrégation à laquelle il avait fait tant de bien, qui lui devait tant de reconnaissance et qui pouvait recueillir encore de si grands avantages de sa présence, eut été bien heureuse de se replacer sous sa conduite. Mais les jugements de Dieu ne sont pas toujours semblables aux jugements des hommes ; le temps où le P. Roger et la Congrégation devaient se retrouver ensemble était encore éloigné ; le Seigneur avait alors d'autres desseins sur nous.
% Source : PDF original, vue 155, côté droit, page manuscrite 311.
Dans sa miséricorde il avait destiné un autre de ses serviteurs à continuer parmi nous le bien que le bon père avait commencé. C'était aux soins de ce nouveau jardinier que Dieu avait réservé le mérite de faire germer, mûrir et fructifier le bon grain dont le P. Roger avait ensemencé le champ auquel il avait donné sa première culture.

Ainsi que nous l'avons dit dans le chapitre précédent, M. Courbon avait donné à M. de Lupé le titre de notre directeur, mais ce n'était qu'un titre, une vraie sinécure. M. Courbon savait tout aussi bien que nous, que malgré tout l'attachement que portait ce vertueux prêtre à une Congrégation dont il faisait lui-même partie, il lui était absolument impossible de remplir les fonctions qu'on lui avait confiées, par la forme. M. le vicaire général cherchait donc sérieusement à nous donner un guide, mais il voulait qu'il nous fût agréable, car il était si bon pour nous qu'il n'aurait pas consenti à nous l'imposer d'autorité ; et il fallait encore que ce guide pût s'occuper effectivement de la Congrégation, la conduire prudemment et sagement, et la diriger dans toutes les bonnes œuvres qu'elle voudrait entreprendre.

Reconnaissons encore ici, Messieurs, comme nous pouvons le faire dans toutes les circonstances un peu importantes de cette histoire, reconnaissons la protection de Marie. Comme elle a soin de nous ! Comme elle veille sur nous ! Comme les prières qu'elle adresse sans cesse à son Divin fils, sont puissantes ! Notre bon Dieu dans sa miséricorde inspira donc en ce moment à M. Courbon une pensée qui était une véritable pensée de salut pour la Congrégation. J'avance ici avec une entière conviction, que c'est à cette Sainte et salutaire pensée que
% Source : PDF original, vue 156, côté gauche, page manuscrite 312.
la Congrégation a dû sa conservation, son accroissement et une grande partie du bien qui s'est opéré par son entremise, et Dieu seul sait ceux d'entre nous que cette pensée a déjà fait entrer dans le ciel et ceux qui grâce à elle pourront y parvenir.

Cette pensée, Messieurs, était celle de jeter les yeux sur M. Recorbet pour lui confier le soin de la Congrégation. La Restauration l'avait rendu au diocèse ; il venait d'être pourvu d'un canonicat à St Jean, ce qui le fixait à Lyon d'une manière définitive.

\section*{M. Recorbet père de la Congrégation}
% Mention marginale de la page manuscrite 312 ; placée avant la description du nouveau père.
Il eut été certainement difficile de rencontrer un homme plus éminemment capable de remplir convenablement les fonctions qu'on allait lui confier. M. Recorbet se présentait à nous précédé par une réputation qui était bien faite pour nous prévenir en sa faveur. Nous étions loin cependant de nous douter du trésor dont la Congrégation allait entrer en possession. Tout en lui excitait à la vertu ; il pouvait bien dire avec le divin maître : apprenez de moi que je suis doux et humble de cœur. Cette douceur évangélique, cette humilité chrétienne qui portaient avec elles un charme ineffable, se manifestaient jusque dans les traits de sa physionomie, par une expression de candeur qui ne se démentait jamais. Il eut suffi aux bons Congréganistes de jeter les yeux sur leur père, pour s'exciter à la pratique de cette vertu d'humilité qui fait les Saints. Il savait répandre sur les paroles d'édification qui sortaient de sa bouche, une onction toute pleine de suavité ; aussi attirait-il à lui tous les cœurs. Il animait le zèle de tous en réchauffant en eux l'amour de Dieu ; il soutenait et consolait les foibles ; rempli comme il l'était de l'esprit de Dieu, il
% Source : PDF original, vue 156, côté droit, page manuscrite 313.
pouvait offrir aux chefs de la Congrégation des lumières certaines pour les empêcher de s'égarer dans les résolutions qu'ils avaient à prendre.

À peine fut-il installé qu'on put commencer à s'apercevoir de tous les avantages qui devaient en résulter. Dès que ce bon père eut adopté sa nouvelle famille, il s'en occupa activement et jusqu'à la fin de ses jours, elle n'a pas cessé d'être l'objet constant de ses sollicitudes. Profitant avec avantage de l'ascendant que lui donnait ce parfum de la bonne odeur de J. C. dont on était vraiment embaumé dès qu'on l'approchait, il s'appliqua surtout à faire naître dans les Congréganistes cet esprit intérieur qui apprend à s'occuper d'abord à se réformer soi-même, afin de pouvoir ensuite travailler avec plus d'efficacité à la gloire de Dieu et à la conversion du prochain.

Disons-le avec franchise, cet esprit manquait peut-être un peu à la Congrégation. Ses membres étaient en général remplis de zèle et d'ardeur pour la défense de la foi et le soutien de la religion ; mais il n'y en avait qu'un bien petit nombre qui fussent assez dégagés de l'esprit du monde pour faire des progrès rapides dans cette vie spirituelle qui dispose les âmes à tous les sacrifices que Dieu pourrait exiger d'elles. Là cependant pouvaient se trouver les moyens de rendre la société véritablement utile, et d'y réchauffer cet esprit de foi qui peut transporter des montagnes.

M. Recorbet chercha à rendre habituelle et journalière pour tous les Congréganistes la pratique de la méditation. Il établit aussi l'usage des retraites. C'est par ces moyens qu'il cherchait à nous bien pénétrer des devoirs d'un véritable Congréganiste,
% Source : PDF original, vue 157, côté gauche, page manuscrite 314.
qui peut se considérer comme une espèce de religieux obligé par état de rester au milieu du monde, mais qui n'en doit pas moins vivre toujours de la vie de la foi.

La grâce que le Seigneur venait d'accorder à la Congrégation, en la mettant entre les mains de M. Recorbet était une nouvelle preuve de sa bienveillance en sa faveur ; mais on peut penser aussi que notre bon maître voulait en lui donnant un pareil guide, la disposer de plus en plus à accomplir les desseins qu'il avait sur elle.

Il nous reste maintenant à suivre l'accomplissement de ces desseins ; mais ici doit s'arrêter cette histoire ; leur développement se trouvera dans les mémoires qui vont suivre.

Déjà en repassant les souvenirs de ces quinze premières années, nous avons eu bien des occasions d'admirer la conduite de la Providence sur nous et de remercier notre bonne mère de sa protection. Nous avons vu successivement notre Congrégation protégée dans son enfance, fortifiée dans sa jeunesse, soutenue et encouragée dans ses épreuves, sauvée presque miraculeusement dans la persécution. Nous venons de la voir restaurée et renouvellée sous les préfectures de MM. du Coin et Garnier. La Providence vient de mettre le comble à ses bienfaits en lui donnant un Saint pour père. Louons donc, bénissons et glorifions le Seigneur en terminant ici cette histoire. Les mémoires que nous allons maintenant entreprendre nous fourniront encore bien des motifs de le louer, de le bénir, de le glorifier de nouveau.

\begin{center}
Fin de l'histoire de la Congrégation des jeunes gens.

\textit{Memor esto, Domine, Congregationis tuæ quam possedisti ab initio.}\footnote{(Note du transcripteur) Littéralement : « Souvenez-vous, Seigneur, de votre assemblée que vous avez acquise dès le commencement ». Psaume 73 (74), verset 2.}
\end{center}
